\documentclass[10pt]{article}

% --- Pacotes Essenciais ---
\usepackage[utf8]{inputenc}
\usepackage[portuguese]{babel}
\usepackage[T1]{fontenc}
\usepackage{graphicx}
\usepackage{amsmath, amssymb}
\usepackage{geometry}
\usepackage{booktabs}
\usepackage{setspace}
\usepackage{float}

% Margens conforme padrão académico
\geometry{a4paper, margin=2.5cm}

\begin{document}

% ======================================================
% PÁGINA DE ROSTO (CAPA)
% ======================================================
\begin{titlepage}
    \centering
    \vspace*{1cm}
    
    
\includegraphics[width=0.4\textwidth]{logo_uc.png}\\[1.5cm]
    \vspace*{2cm}
    
    {\large Departamento de Engenharia Informática} \\[1.5cm]
    
    \rule{\textwidth}{1pt} \\[0.4cm]
    {\huge \bfseries Inteligência Artificial para Entretenimento} \\[0.2cm]
    {\large \bfseries Projeto 1: Co-Creative Curation Pipeline on MovieLens 100K} \\[0.4cm]
    \rule{\textwidth}{1pt} \\[1.5cm]
    
    \begin{minipage}{0.8\textwidth}
        \centering
        \large
        \textbf{Elementos do grupo:} \\[0.8cm]
        Miguel Soares Francisco -- Nº 2025179328 -- uc2025179328@student.uc.pt \\
        Miguel Oliveira Pinto -- Nº 2025190104 -- uc2025190104@student.uc.pt \\
        Gonçalo José Carrajola Gaio -- Nº 2022224905 -- uc2022224905@student.uc.pt \\
    \end{minipage}

    \vfill
    
    {\large Coimbra, 2025-2026}
\end{titlepage}

% ======================================================
% ÍNDICE
% ======================================================
\newpage
\onehalfspacing
\tableofcontents
\newpage

% ======================================================
% SECÇÕES DO RELATÓRIO
% ======================================================

\section{Introdução}

No contexto da disciplina de Inteligência Artificial para Entretenimento, os sistemas de recomendação têm um papel fundamental na seleção e organização de conteúdos relevantes para cada utilizador. Neste projeto, esse problema é abordado sob a perspetiva de \textit{co-creative curation}, em que a recomendação é vista como um processo iterativo que combina geração de candidatos, ranking, reranking, feedback e atualização da personalização. Com base no dataset MovieLens 100K, foi desenvolvido um pipeline completo capaz de produzir listas top-$K$ de filmes adaptadas às preferências dos utilizadores. Para este efeito, foram implementados e comparados diferentes métodos de ranking, incluindo um baseline de popularidade, modelos de fatorização matricial com SGD e ALS, e um modelo de \textit{Pairwise Learning-to-Rank}. Adicionalmente, foi incorporado um mecanismo de reranking com MMR para controlar a diversidade das recomendações e um mecanismo de atualização do estado do utilizador com EMA para simular sessões sucessivas de interação. A avaliação experimental foi realizada com recurso às métricas Recall@10, NDCG@10 e Diversity@10, permitindo analisar não só a relevância das recomendações, mas também o equilíbrio entre precisão, diversidade e adaptação dinâmica ao utilizador.

\section{Arquitetura do Sistema}

O sistema segue a sequência exigida no enunciado: geração de candidatos, ranking, reranking com controlo de diversidade, feedback e atualização do estado do utilizador. A Tabela~\ref{tab:estrutura_ficheiros} resume os principais módulos que implementam cada bloco do pipeline e a respetiva função no fluxo de execução.

\begin{table}[H]
\centering
\small
\renewcommand{\arraystretch}{1.35}
\begin{tabular}{p{0.34\textwidth} p{0.58\textwidth}}
\toprule
\textbf{Ficheiro / Módulo} & \textbf{Função} \\
\midrule
\texttt{src/main.py} & Orquestra toda a execução do projeto: configuração, treino dos modelos, avaliação, aplicação de MMR, sessões com EMA e criação de logs. \\
\texttt{src/core/config.py} & Define os principais parâmetros experimentais, hiperparâmetros dos modelos e opções de execução, incluindo o modo FAST MODE. \\
\texttt{src/core/structs.py} & Contém as estruturas de dados e classes auxiliares usadas para guardar resultados, variáveis globais e modelos. \\
\texttt{src/core/candidates.py} & Gera o conjunto de filmes candidatos para cada utilizador, excluindo os itens já observados no treino. \\
\texttt{src/core/utils.py} & Reúne funções auxiliares para carregamento de dados, preparação da avaliação, geração de gráficos, comparação de métodos e escrita de logs. \\
\texttt{src/core/metrics.py} & Implementa as métricas de avaliação usadas no projeto, nomeadamente Recall@10, NDCG@10, similaridade do cosseno e Diversity@10. \\
\texttt{src/rankers/popularity.py} & Implementa o baseline de popularidade, responsável por recomendar filmes com base no número de interações no treino. \\
\texttt{src/rankers/mf\_general.py} & Contém funções gerais para os modelos de fatorização matricial, incluindo preparação dos dados, previsão de ratings, recomendação e cálculo de RMSE. \\
\texttt{src/rankers/mf\_sgd.py} & Implementa o treino do modelo de Matrix Factorization com SGD, incluindo fatores latentes, biases e histórico de treino. \\
\texttt{src/rankers/mf\_als.py} & Implementa o treino do modelo de Matrix Factorization com ALS, atualizando alternadamente fatores latentes e biases. \\
\texttt{src/rankers/pairwise\_ltr.py} & Implementa o modelo Pairwise Learning-to-Rank, incluindo criação dos pares de treino, construção das features, treino e recomendação. \\
\texttt{src/reranker/mmr.py} & Implementa o reranking com MMR para aumentar a diversidade das listas recomendadas, equilibrando relevância e diversidade. \\
\texttt{src/personalization/ema.py} & Implementa a atualização dinâmica do estado do utilizador com EMA e a simulação de sessões iterativas de recomendação. \\
\bottomrule
\end{tabular}
\caption{Descrição resumida dos principais ficheiros do projeto.}
\label{tab:estrutura_ficheiros}
\end{table}

\section{Conjunto de dados e pré-processamento}

O sistema foi desenvolvido com base no \textit{MovieLens 100K}, com avaliações explícitas e metadados de filmes (título e géneros). Estes géneros são essenciais para a medida de diversidade e para a similaridade no MMR. A Tabela~\ref{tab:dataset_preproc} sintetiza as decisões principais e os números usados no pré-processamento.

\begin{table}[H]
\centering
\small
\renewcommand{\arraystretch}{1.2}
\begin{tabular}{p{0.28\textwidth} p{0.62\textwidth}}
\toprule
\textbf{Aspeto} & \textbf{Decisão/Resumo} \\
\midrule
Dataset & MovieLens 100K; 100\,000 avaliações, 943 utilizadores, 1682 filmes; tabela de itens com 24 colunas. \\
Split & Divisão aleatória 80/20: 80\,000 interações em treino e 20\,000 em teste. \\
Relevância & $G_u = \{ i : (u,i)\in\Omega_{\text{test}} \wedge r_{ui} \ge 4 \}$; itens do treino excluídos das recomendações. \\
Utilizadores elegíveis & 920 com pelo menos um item relevante; avaliação em 300 utilizadores (modo rápido). \\
Representação MF & Mapeamento para índices consecutivos; dados em tuplos $(u,i,r)$; estruturas por utilizador/item para ALS. \\
LTR pairwise & 46\,409 pares de preferência $(i^+, i^-)$ gerados a partir do treino. \\
Diversidade/MMR & Vetores binários de géneros; similaridade por cosseno. \\
\bottomrule
\end{tabular}
\caption{Resumo do conjunto de dados e pré-processamento.}
\label{tab:dataset_preproc}
\end{table}

Em termos operacionais, as avaliações foram carregadas e divididas em treino e teste, aplicando a definição de relevância do enunciado e excluindo itens já observados no treino durante a recomendação. Para os modelos de fatorização matricial, a indexação consecutiva permitiu construir as matrizes latentes e as estruturas auxiliares necessárias ao SGD e ao ALS. Para o \textit{pairwise LTR}, geraram-se pares de preferência por utilizador, alinhando o treino com a tarefa de ordenação. Finalmente, os vetores de género foram preparados para suportar tanto o MMR como o cálculo de Diversity@10.






\section{Implementação do pipeline}

O fluxo completo de \textit{co-creative curation} está concentrado em \texttt{main.py} e segue o percurso exigido: geração de candidatos, ranking, reranking, feedback e atualização do estado do utilizador. O pipeline avalia primeiro o baseline de popularidade, treina e avalia os modelos de fatorização matricial e o \textit{pairwise LTR}, aplica MMR sobre os top-\(M\) candidatos e, por fim, simula sessões com EMA.

A Tabela~\ref{tab:pipeline_impl} sintetiza como cada etapa foi concretizada. Em todas as fases, os resultados são registados em JSONL com identificador interno do utilizador, top-\(K\), métricas e hiperparâmetros; nas sessões EMA, inclui-se o item escolhido e um resumo do estado atualizado.

\begin{table}[H]
\centering
\small
\renewcommand{\arraystretch}{1.2}
\begin{tabular}{p{0.22\textwidth} p{0.68\textwidth}}
\toprule
\textbf{Etapa} & \textbf{Implementação (\texttt{main.py})} \\
\midrule
Candidatos & \(C(u)\) inclui todos os itens não avaliados no treino; usado por todos os rankers. \\
Popularidade & Avaliação do baseline com ordenação por contagem global. \\
MF-SGD / MF-ALS & Treino com vieses e avaliação com Recall@10, NDCG@10 e Diversity@10. \\
Pairwise LTR & Pares \((i^+, i^-)\) gerados no treino; modelo Bradley--Terry com \(f(u,i)=[s_{MF}, b_u, b_i, pop(i), 1]\). \\
MMR & Reranking dos top-\(M\) com \(\alpha \in \{0.1,0.4,0.7\}\) e similaridade por cosseno dos géneros. \\
EMA & Sessões com 5 rondas, escolha heurística e atualização \(u_{t+1} = \frac{(1-\rho)u_t + \rho v_i}{\|(1-\rho)u_t + \rho v_i\|}\), usando \(\rho\in\{0.1,0.3\}\). \\
Logs & JSONL por método e por sessão, com top-\(K\), métricas e hiperparâmetros. \\
\bottomrule
\end{tabular}
\caption{Resumo da implementação do pipeline em \texttt{main.py}.}
\label{tab:pipeline_impl}
\end{table}

O modo rápido controla o número de utilizadores avaliados e os principais hiperparâmetros (dimensão latente, épocas/iterações e \(M\) do MMR), garantindo execução completa em tempo reduzido, conforme requerido.


\subsection{Configuração experimental}
Para clarificar as definições usadas, a Tabela~\ref{tab:config_execucao} resume os parâmetros mais relevantes no modo rápido e no modo completo. Os restantes hiperparâmetros (por exemplo, \(K=10\), \(\alpha\) do MMR e \(\rho\) do EMA) são mantidos constantes entre modos, conforme indicado no enunciado.


\begin{table}[H]
\centering
\small
\renewcommand{\arraystretch}{1.2}
\begin{tabular}{p{0.36\textwidth} p{0.25\textwidth} p{0.25\textwidth}}
\toprule
\textbf{Parâmetro} & \textbf{Modo rápido} & \textbf{Modo completo} \\
\midrule
Utilizadores avaliados & 300 & Todos os elegíveis \\
Dimensão latente (MF) & 20 & 50 \\
Épocas MF-SGD & 10 & 25 \\
Iterações MF-ALS & 8 & 15 \\
Top-\(M\) para MMR & 80 & 200 \\
\midrule
Top-\(K\) & 10 & 10 \\
\(\alpha\) do MMR & \(\{0.1,0.4,0.7\}\) & \(\{0.1,0.4,0.7\}\) \\
\(\rho\) do EMA & \(\{0.1,0.3\}\) & \(\{0.1,0.3\}\) \\
Pares por utilizador (LTR) & 50 & 50 \\
Épocas LTR & 5 & 5 \\
\bottomrule
\end{tabular}
\caption{Resumo das configurações de execução.}
\label{tab:config_execucao}
\end{table}
\noindent\textit{Nota:} o modo rápido é controlado pela \textit{flag} \texttt{FAST\_MODE} no ficheiro \texttt{config.py}.

















\section{Resultados}

Nesta secção apresentam-se os principais resultados obtidos para os diferentes métodos implementados no pipeline de recomendação, considerando as métricas Recall@10, NDCG@10 e Diversity@10. Os valores reportados correspondem ao modo completo (\texttt{FAST\_MODE} desativado), com avaliação sobre todos os utilizadores elegíveis.

\subsection{Comparação entre métodos base}

A Tabela~\ref{tab:main-results} resume os resultados médios obtidos pelos quatro métodos principais avaliados: baseline de popularidade, fatorização matricial com SGD, fatorização matricial com ALS e modelo \textit{pairwise learning-to-rank} (E1). No modo completo, o baseline de popularidade apresenta os melhores valores de Recall@10 (0.1372) e NDCG@10 (0.1764), o que indica que a popularidade global continua a ser um sinal forte neste dataset.

O MF-SGD melhora relativamente ao MF-ALS, mas continua abaixo do baseline em relevância; o MF-ALS mantém o pior desempenho em Recall@10 e NDCG@10. O \textit{pairwise LTR} apresentou valores inferiores aos do baseline (Recall@10 de 0.0314 e NDCG@10 de 0.0470), mantendo diversidade semelhante. Este comportamento reforça a necessidade de analisar o impacto das features e das ablações (Secção de ablação).

Uma interpretação plausível para o desempenho do LTR é a escala heterogénea das features (por exemplo, a popularidade em contagens absolutas) e o facto de a previsão MF estar limitada ao intervalo de ratings. Sem normalização explícita, o modelo linear pode favorecer sinais dominantes, o que ajuda a explicar a melhoria observada quando a popularidade é removida.


\begin{table}[H]
\centering
\caption{Comparação dos métodos base na avaliação offline.}
\label{tab:main-results}
\begin{tabular}{lccc}
\toprule
\textbf{Método} & \textbf{Recall@10} & \textbf{NDCG@10} & \textbf{Diversity@10} \\
\midrule
Popularidade & 0.1372 & 0.1764 & 0.7517 \\
MF-SGD & 0.0506 & 0.0805 & 0.7338 \\
MF-ALS & 0.0196 & 0.0297 & 0.7086 \\
Pairwise LTR & 0.0314 & 0.0470 & 0.7511 \\
\bottomrule
\end{tabular}
\end{table}

De forma geral, os resultados indicam que otimizar diretamente para ordenação relativa entre itens foi mais eficaz do que usar exclusivamente modelos de fatorização matricial treinados para previsão de ratings. O desempenho do modelo pairwise LTR confirma, assim, a importância de alinhar o objetivo de treino com a tarefa final de ranking. 

\subsection{Impacto do reranking com MMR}

Um dos requisitos experimentais do projeto (E2) consiste em apresentar um gráfico do compromisso entre relevância e diversidade, medido em termos de NDCG@10 e Diversity@10, para diferentes valores do parâmetro \(\alpha \in \{0.1, 0.4, 0.7\}\), aplicando o método MMR pelo menos a um baseline de fatorização matricial. Para satisfazer esse requisito, foi analisado o efeito do reranking com MMR sobre o modelo MF-SGD, representando graficamente a relação entre NDCG@10 e Diversity@10 para os diferentes valores de \(\alpha\).

Após a avaliação dos rankers base, o método MMR foi aplicado para introduzir controlo explícito da diversidade nas listas finais. A Tabela~\ref{tab:mmr-results} apresenta os resultados médios para diferentes valores de \(\alpha\), usando como base os modelos MF-SGD, MF-ALS e Pairwise LTR, enquanto a Figura~\ref{fig:mmr_tradeoff_mf_sgd} ilustra graficamente o trade-off obtido no caso do MF-SGD. 

No caso do MF-SGD, o reranking com MMR aumentou significativamente a diversidade: de 0.7338 sem reranking para 0.7957, 0.9304 e 0.9909 para \(\alpha = 0.1\), \(0.4\) e \(0.7\), respetivamente. Em paralelo, os valores de NDCG@10 passaram de 0.0805 (sem MMR) para 0.0790, 0.0718 e 0.0579 à medida que a diversidade aumentou. A Figura~\ref{fig:mmr_tradeoff_mf_sgd} ilustra este compromisso, com \(\alpha = 0.4\) a oferecer um equilíbrio entre diversidade elevada e perda moderada de relevância.

Para o MF-ALS, o MMR também aumentou muito a diversidade, elevando-a para valores próximos de 0.96. No entanto, como o modelo base já apresentava fraca qualidade de ranking, o reranking não foi suficiente para recuperar relevância. Já no caso do Pairwise LTR, observou-se um comportamento mais equilibrado: a diversidade aumentou progressivamente com \(\alpha\), enquanto Recall@10 e NDCG@10 sofreram apenas uma ligeira degradação para valores mais elevados deste parâmetro. A Figura~\ref{fig:mmr_tradeoff_outros} mostra estes trade-offs adicionais e reforça o compromisso clássico entre relevância e diversidade.


\begin{figure}[H]
    \centering
    \includegraphics[width=0.72\textwidth]{images/mmr_tradeoff_mf_sgd.png}
    \caption{Trade-off entre NDCG@10 e Diversity@10 para o reranking com MMR aplicado ao modelo MF-SGD, considerando \(\alpha \in \{0.1, 0.4, 0.7\}\).}
    \label{fig:mmr_tradeoff_mf_sgd}
\end{figure}

\begin{figure}[H]
    \centering
    \begin{minipage}{0.48\textwidth}
        \centering
        \includegraphics[width=\textwidth]{images/mmr_tradeoff_mf_als.png}
    \end{minipage}\hfill
    \begin{minipage}{0.48\textwidth}
        \centering
        \includegraphics[width=\textwidth]{images/mmr_tradeoff_pairwise_ltr.png}
    \end{minipage}
    \caption{Trade-off do MMR nos modelos MF-ALS e Pairwise LTR (modo completo).}
    \label{fig:mmr_tradeoff_outros}
\end{figure}


\begin{table}[H]
\centering
\caption{Resultados do reranking com MMR para diferentes valores de \(\alpha\).}
\label{tab:mmr-results}
\begin{tabular}{llccc}
\toprule
\textbf{Modelo Base} & \textbf{\(\alpha\)} & \textbf{Recall@10} & \textbf{NDCG@10} & \textbf{Diversity@10} \\
\midrule
MF-SGD & 0.1 & 0.0492 & 0.0790 & 0.7957 \\
MF-SGD & 0.4 & 0.0444 & 0.0718 & 0.9304 \\
MF-SGD & 0.7 & 0.0295 & 0.0579 & 0.9909 \\
\midrule
MF-ALS & 0.1 & 0.0120 & 0.0230 & 0.9751 \\
MF-ALS & 0.4 & 0.0120 & 0.0230 & 0.9832 \\
MF-ALS & 0.7 & 0.0123 & 0.0234 & 0.9916 \\
\midrule
Pairwise LTR & 0.1 & 0.0302 & 0.0459 & 0.8066 \\
Pairwise LTR & 0.4 & 0.0281 & 0.0424 & 0.9459 \\
Pairwise LTR & 0.7 & 0.0191 & 0.0350 & 0.9949 \\
\bottomrule
\end{tabular}
\end{table}

Entre os resultados obtidos, destaca-se particularmente o Pairwise LTR com MMR e \(\alpha = 0.4\), que preserva um NDCG@10 muito próximo do melhor valor absoluto observado, ao mesmo tempo que aumenta claramente a diversidade da lista. Este comportamento mostra que a fase de reranking pode acrescentar valor mesmo quando o ranker base já apresenta boa qualidade.

\subsection{Sessões de personalização com EMA}

Para além da avaliação estática, foi ainda demonstrada a componente dinâmica de personalização através de sessões sucessivas com atualização por EMA. Foram realizadas sessões com 3 utilizadores, 5 rondas por sessão e comparação entre \(\rho = 0.1\) e \(\rho = 0.3\). Em cada ronda, o sistema produziu uma lista recomendada, simulou a escolha de um item e atualizou o vetor de estado do utilizador com base no item selecionado.

Os resultados qualitativos mostram alterações graduais das listas ao longo das rondas. No MF-SGD, a sobreposição média entre listas consecutivas manteve-se elevada para \(\rho = 0.1\) (entre 0.768 e 0.818), enquanto com \(\rho = 0.3\) desceu para intervalos entre 0.624 e 0.725, indicando maior sensibilidade ao feedback recente. No MF-ALS, observou-se padrão semelhante: \(\rho = 0.1\) oscilou entre 0.674 e 0.818, e \(\rho = 0.3\) entre 0.465 e 0.632. Estas tendências são coerentes com a interpretação de \(\rho\) como taxa de atualização.

\begin{figure}[H]
    \centering
    \begin{minipage}{0.48\textwidth}
        \centering
        \includegraphics[width=\textwidth]{images/ema_rho_overlap_mf_sgd.png}
    \end{minipage}\hfill
    \begin{minipage}{0.48\textwidth}
        \centering
        \includegraphics[width=\textwidth]{images/ema_rho_overlap_mf_als.png}
    \end{minipage}
    \caption{Efeito de \(\rho\) na sobreposição média entre listas consecutivas (Jaccard) em sessões EMA com MF-SGD e MF-ALS.}
    \label{fig:ema_rho_overlap}
\end{figure}

Estes resultados confirmam que o mecanismo de EMA implementado é capaz de introduzir adaptação temporal no pipeline, permitindo que o sistema reaja ao feedback do utilizador ao longo de uma sessão, em vez de produzir sempre listas estáticas.

\subsection{Curvas de treino (modo completo)}
Para complementar a análise, apresentam-se as curvas de treino dos modelos MF e LTR no modo completo. Observa-se uma descida consistente do RMSE nos modelos MF e uma estabilização da perda no LTR, indicando convergência do processo de treino.

\begin{figure}[H]
    \centering
    \begin{minipage}{0.48\textwidth}
        \centering
        \includegraphics[width=\textwidth]{images/mf_sgd_training_rmse.png}
    \end{minipage}\hfill
    \begin{minipage}{0.48\textwidth}
        \centering
        \includegraphics[width=\textwidth]{images/mf_als_training_rmse.png}
    \end{minipage}

    \vspace{0.2cm}
    \begin{minipage}{0.62\textwidth}
        \centering
        \includegraphics[width=\textwidth]{images/pairwise_ltr_training_loss.png}
    \end{minipage}
    \caption{Curvas de treino no modo completo: RMSE do MF-SGD e MF-ALS e perda do Pairwise LTR.}
    \label{fig:training_curves_full}
\end{figure}





\section{Testes adicionais}

De acordo com o enunciado do projeto (E3), foram considerados estudos com testes adicionais com o objetivo de perceber o contributo de componentes específicas do pipeline, nomeadamente o mecanismo de reranking com MMR e a utilização da feature de popularidade no modelo \textit{pairwise learning-to-rank}.

\subsection{Remoção do MMR}

Comparamos o desempenho dos modelos sem reranking e com MMR para dois níveis de \(\alpha\). O padrão é claro: o MMR aumenta substancialmente a diversidade, mas reduz a relevância (Recall@10 e NDCG@10), especialmente para \(\alpha=0.7\).

\begin{table}[H]
\centering
\small
\renewcommand{\arraystretch}{1.2}
\begin{tabular}{llccc}
\toprule
\textbf{Modelo} & \textbf{Configuração} & \textbf{Recall@10} & \textbf{NDCG@10} & \textbf{Diversity@10} \\
\midrule
MF-SGD & Sem MMR & 0.0506 & 0.0805 & 0.7338 \\
MF-SGD & MMR (\(\alpha=0.4\)) & 0.0444 & 0.0718 & 0.9304 \\
MF-SGD & MMR (\(\alpha=0.7\)) & 0.0295 & 0.0579 & 0.9909 \\
\midrule
Pairwise LTR & Sem MMR & 0.0314 & 0.0470 & 0.7511 \\
Pairwise LTR & MMR (\(\alpha=0.4\)) & 0.0281 & 0.0424 & 0.9459 \\
Pairwise LTR & MMR (\(\alpha=0.7\)) & 0.0191 & 0.0350 & 0.9949 \\
\bottomrule
\end{tabular}
\caption{Impacto do MMR nas métricas de ranking e diversidade.}
\label{tab:ablation_mmr}
\end{table}

\subsection{Remoção da feature de popularidade no Pairwise LTR}

Este teste compara o modelo pairwise LTR com e sem a feature de popularidade. Os resultados mostram um aumento de relevância ao remover a popularidade, acompanhado de uma ligeira redução da diversidade.

\begin{table}[H]
\centering
\small
\renewcommand{\arraystretch}{1.2}
\begin{tabular}{lccc}
\toprule
\textbf{Configuração} & \textbf{Recall@10} & \textbf{NDCG@10} & \textbf{Diversity@10} \\
\midrule
Com popularidade & 0.0314 & 0.0470 & 0.7511 \\
Sem popularidade & 0.0520 & 0.0857 & 0.7219 \\
\bottomrule
\end{tabular}
\caption{Ablation do Pairwise LTR: efeito da feature de popularidade.}
\label{tab:ablation_popularity}
\end{table}











\section{Conclusão}
Este trabalho implementou um pipeline completo de 	extit{co-creative curation} e permitiu comparar baselines simples com modelos mais avançados, avaliando simultaneamente relevância, diversidade e adaptação temporal. Os resultados mostram que a popularidade continua a ser um sinal forte neste dataset, que o reranking por MMR melhora claramente a diversidade com um custo controlável em relevância, e que o EMA introduz uma dinâmica coerente de atualização do perfil do utilizador. As ablações confirmam ainda que o desempenho do 	extit{pairwise LTR} é sensível às features usadas, reforçando a importância de alinhar a representação e o objetivo de treino com a tarefa de ranking.

Quanto à escala das métricas, valores de Recall@10 e NDCG@10 baixos são esperados neste cenário: o catálogo é grande, o corte é restrito a $K=10$ e a relevância é definida por um limiar de rating, o que aumenta o número de itens relevantes por utilizador e comprime as métricas. Assim, apesar de estarem em $[0,1]$, os valores absolutos não são diretamente comparáveis a outros contextos. O desempenho particularmente baixo do MF-ALS é plausível por estar otimizado para erro quadrático em ratings (não para ranking) e por ser sensível a hiperparâmetros e regularização; com as configurações usadas, pode subajustar e produzir ordenações menos discriminativas do que o MF-SGD. No conjunto, o projeto evidencia o compromisso clássico entre precisão e variedade e sublinha que a qualidade final depende tanto do ranker base como das etapas de reranking e personalização.

\end{document}

% \newpage

% % ======================================================
% % APÊNDICE TEÓRICO (OBRIGATÓRIO T1-T5)
% % ======================================================
% \appendix
% \section{Apêndice Teórico: Respostas T1-T5}


% 
\section*{T1 (ALS closed-form update)}

\subsection*{1. Redução a um Problema de Regressão Ridge}

O objetivo da Fatoração de Matrizes (MF) é minimizar o erro quadrático entre a classificação real $r_{ui}$ e a prevista $\hat{r}_{ui}$. A previsão é definida como:
\begin{equation}
\hat{r}_{ui} = \mu + b_u + b_i + p_u^\top q_i
\end{equation}

Para isolar a otimização do vetor do utilizador $p_u$, definimos o \textbf{resíduo} $e_{ui}$ como a parte do rating que não é explicada pelos termos fixos ($\mu, b_u, b_i$):
\begin{equation}
e_{ui} = r_{ui} - (\mu + b_u + b_i)
\end{equation}

Ao fixar os vetores dos itens $Q$, a função de custo local para um utilizador $u$ com regularização $L_2$ torna-se:
\begin{equation}
\min_{p_u} \sum_{i \in \Omega(u)} (e_{ui} - p_u^\top q_i)^2 + \lambda \|p_u\|^2
\end{equation}

Esta expressão é a forma canónica de uma \textbf{Regressão Ridge}, onde $p_u$ atua como o vetor de pesos, $q_i$ como as características (features) e $e_{ui}$ como o valor alvo.

\subsection*{2. Derivação das Equações Normais}

Para encontrar o valor ótimo de $p_u$, calculamos o gradiente da função de custo $J(p_u)$ e igualamos a zero:

\textbf{Passo A: Gradiente}
\begin{equation}
\nabla_{p_u} J(p_u) = \sum_{i \in \Omega(u)} -2(e_{ui} - p_u^\top q_i)q_i + 2\lambda p_u
\end{equation}

\textbf{Passo B: Condição de Primeira Ordem}
Igualando a zero e simplificando o fator 2:
\begin{equation}
\sum_{i \in \Omega(u)} (p_u^\top q_i)q_i + \lambda p_u = \sum_{i \in \Omega(u)} e_{ui} q_i
\end{equation}

\textbf{Passo C: Formulação Matricial}
Reescrevendo o produto escalar $(p_u^\top q_i)q_i$ como $(q_i q_i^\top) p_u$ e utilizando a matriz Identidade $I$:
\begin{equation}
\left( \sum_{i \in \Omega(u)} q_i q_i^\top + \lambda I \right) p_u = \sum_{i \in \Omega(u)} e_{ui} q_i
\end{equation}

\textbf{Passo D: Solução de Forma Fechada}
Multiplicando pelo inverso da matriz resultante, obtemos a regra de atualização do ALS:
\begin{equation}
p_u = \left( \sum_{i \in \Omega(u)} q_i q_i^\top + \lambda I \right)^{-1} \left( \sum_{i \in \Omega(u)} e_{ui} q_i \right)
\end{equation}

\section*{T2 (SGD updates for MF with biases)}

A função de previsão para a classificação $\hat{r}_{ui}$ é definida como:
\begin{equation}
\hat{r}_{ui} = \mu + b_u + b_i + p_u^\top q_i
\end{equation}

O objetivo é minimizar o erro quadrático com regularização $\ell_2$ para cada observação $(u, i) \in \Omega$. A função de custo local $J_{ui}$ é:
\begin{equation}
J_{ui} = (r_{ui} - \hat{r}_{ui})^2 + \lambda (\|p_u\|^2 + \|q_i\|^2 + b_u^2 + b_i^2)
\end{equation}

Definimos o erro de previsão como $\epsilon_{ui} = r_{ui} - \hat{r}_{ui}$. Para derivar as atualizações SGD, calculamos o gradiente negativo de $J_{ui}$ em ordem a cada parâmetro.

\subsection*{1. Derivação dos Gradientes}
Aplicando a regra da cadeia, os gradientes em relação aos parâmetros são:
\begin{itemize}
    \item $\frac{\partial J_{ui}}{\partial p_u} = -2(r_{ui} - \hat{r}_{ui})q_i + 2\lambda p_u = -2\epsilon_{ui}q_i + 2\lambda p_u$
    \item $\frac{\partial J_{ui}}{\partial q_i} = -2(r_{ui} - \hat{r}_{ui})p_u + 2\lambda q_i = -2\epsilon_{ui}p_u + 2\lambda q_i$
    \item $\frac{\partial J_{ui}}{\partial b_u} = -2(r_{ui} - \hat{r}_{ui}) + 2\lambda b_u = -2\epsilon_{ui} + 2\lambda b_u$
    \item $\frac{\partial J_{ui}}{\partial b_i} = -2(r_{ui} - \hat{r}_{ui}) + 2\lambda b_i = -2\epsilon_{ui} + 2\lambda b_i$
\end{itemize}

\subsection*{2. Regras de Atualização}
Seguindo a direção oposta ao gradiente ($\theta \leftarrow \theta - \eta \nabla J$) e absorvendo o fator 2 na taxa de aprendizagem $\eta$, obtemos as seguintes regras de atualização:
\begin{itemize}
    \item $p_u \leftarrow p_u + \eta (\epsilon_{ui} q_i - \lambda p_u)$
    \item $q_i \leftarrow q_i + \eta (\epsilon_{ui} p_u - \lambda q_i)$
    \item $b_u \leftarrow b_u + \eta (\epsilon_{ui} - \lambda b_u)$
    \item $b_i \leftarrow b_i + \eta (\epsilon_{ui} - \lambda b_i)$
\end{itemize}





\section*{T3 (Pairwise gradient)}

O objetivo é demonstrar que o gradiente da função de perda logística em relação aos parâmetros $\phi$ é dado por:
\begin{equation}
\nabla_{\phi} \left( -\log \sigma(\phi^\top \Delta f) \right) = \left( \sigma(\phi^\top \Delta f) - 1 \right) \Delta f
\end{equation}

\subsection*{Demonstração Passo a Passo}

Seja $z = \phi^\top \Delta f$. A função que queremos derivar é $L(\phi) = -\log \sigma(z)$.

\textbf{1. Derivada da função sigmóide ($\sigma$):}
Sabemos que a derivada da função sigmóide $\sigma(z) = \frac{1}{1 + e^{-z}}$ em relação a $z$ é:
\begin{equation}
\sigma'(z) = \sigma(z)(1 - \sigma(z))
\end{equation}

\textbf{2. Aplicação da Regra da Cadeia:}
Utilizando a regra da cadeia para o logaritmo e para a sigmóide:
\begin{equation}
\nabla_{\phi} L(\phi) = -\frac{1}{\sigma(z)} \cdot \frac{\partial \sigma(z)}{\partial z} \cdot \nabla_{\phi} z
\end{equation}

\textbf{3. Substituição das derivadas:}
Substituindo $\sigma'(z)$ e notando que $\nabla_{\phi} (\phi^\top \Delta f) = \Delta f$:
\begin{equation}
\nabla_{\phi} L(\phi) = -\frac{1}{\sigma(z)} \cdot \left[ \sigma(z)(1 - \sigma(z)) \right] \cdot \Delta f
\end{equation}

\textbf{4. Simplificação:}
Os termos $\sigma(z)$ no numerador e denominador anulam-se:
\begin{equation}
\nabla_{\phi} L(\phi) = -(1 - \sigma(z)) \Delta f
\end{equation}
\begin{equation}
\nabla_{\phi} L(\phi) = (\sigma(z) - 1) \Delta f
\end{equation}

\textbf{5. Resultado Final:}
Substituindo $z$ de volta pela expressão original $\phi^\top \Delta f$:
\begin{equation}
\nabla_{\phi} \left( -\log \sigma(\phi^\top \Delta f) \right) = \left( \sigma(\phi^\top \Delta f) - 1 \right) \Delta f
\end{equation}

\section*{T4 (Why RMSE $\neq$ ranking quality?)}

Otimizar o RMSE (Root Mean Squared Error) em modelos de Fatoração de Matrizes nem sempre garante uma boa qualidade de recomendação Top-K devido aos seguintes fatores:

\begin{itemize}
    \item \textbf{Foco na Regressão vs. Ordenação:} O RMSE trata o problema como uma regressão pontual, tentando prever o valor exato de cada classificação. No entanto, para recomendações, o que importa é a ordem relativa entre os itens (ranking) e não a precisão do valor absoluto da nota.
    
    \item \textbf{Peso Igual a Erros Irrelevantes:} O RMSE penaliza igualmente um erro de previsão num item que o utilizador detesta (ex: prever 2 em vez de 1) e num item que o utilizador adoraria (ex: prever 4 em vez de 5). Para o cenário de Top-K, apenas os erros ocorridos no topo da lista são realmente críticos para a experiência do utilizador.
    
    \item \textbf{Desalinhamento com Métricas de Negócio:} O objetivo de um sistema de entretenimento é frequentemente maximizar métricas como NDCG ou Recall, que dependem da posição dos itens relevantes na lista. Um modelo pode reduzir o seu RMSE global ao prever melhor as notas baixas da "cauda longa" (long-tail) sem nunca melhorar a visibilidade dos itens que o utilizador realmente escolheria consumir.
    
    \item \textbf{Viés de Dados Observados:} O RMSE é calculado apenas sobre interações observadas (ratings explícitos). Isto ignora a estrutura de interações implícitas, onde o ato de escolher um item entre várias opções num conjunto de candidatos é um sinal de preferência muito mais forte para o ranking do que a precisão de uma nota isolada.
\end{itemize}



\subsection*{Observações e Resultados do Projeto}
No âmbito do projeto desenvolvido, observou-se que a redução do RMSE durante o treino do MF-ALS não se traduziu numa melhoria equivalente das métricas de ranking. Em particular, apesar de o RMSE descer de forma consistente, o NDCG@10 médio manteve-se baixo (0.0297), enquanto o modelo pairwise LTR, que otimiza diretamente a ordenação, atingiu NDCG@10 de 0.0470. Este contraste confirma que otimizar o erro de previsão de ratings não garante um bom ordenamento dos itens no topo da lista.
\section*{T5 (Offline evaluation caveat)}

A avaliação offline pode ser enganosa porque os dados históricos não são neutros, mas sim o resultado de uma política de exposição anterior que gera um \textbf{viés de seleção}. Itens que nunca foram mostrados constituem "pontos cegos", impossibilitando testar a reação do utilizador a esses conteúdos. Além disso, o modelo pode limitar-se a imitar os enviesamentos do sistema antigo, ignorando a \textbf{natureza de ciclo fechado} dos recomendadores, onde o ranking atua como uma intervenção que molda o comportamento e os dados futuros. Sem exposição aleatória, criam-se \textbf{ciclos de feedback} que amplificam a popularidade e reduzem a diversidade, algo que métricas estáticas como o NDCG não conseguem capturar totalmente por ignorarem a deriva de dados (data drift) causada pelo próprio sistema.




